\section{Effective calculus needed for curved cells}\label{sec:calculus}

The cube theorem is complete without this section. To apply it to a curved
surface, however, the pulled-back coefficients must themselves have the
required certificates. We now supply the relevant constructions and
proofs. In particular, the chain rule will not be assumed as an
unexplained analytic oracle.

\subsection{Full first-order jets}

Let $D\subseteq\Q^d$ be a closed rational box. A certified $C^1$ presentation
consists of coefficient and derivative evaluators $H,J_1,\ldots,J_d$,
coherence and uniform continuity certificates, integer bounds $A,D_1$
on the named value and each named derivative, and a primitive recursive
remainder modulus $\rho$. Explicitly, the bounds mean
$|H(x,p)|\le A+\epsilon_p$ and $|J_j(x,p)|\le D_1+\epsilon_p$.
For rational $x,y$ in the box, put $h=\|y-x\|_\infty$. The full
remainder condition is
\begin{align}
 h\le2^{-\rho(t)}\quad\Longrightarrow\quad&
 \left|H(y,p)-H(x,p)
             -\sum_{j=1}^dJ_j(x,p)(y_j-x_j)\right|\notag\\
 &\le h2^{-t}+(2+dh)2^{-p}.
 \label{eq:full-jet}
\end{align}
For a coordinate increment only one derivative appears. To derive (D)
rigorously from \eqref{eq:full-jet}, let $R_p$ be that coordinate residual.
Coherence gives
$|R_p-R_q|\le(2+h)(\epsilon_p+\epsilon_q)$.
Applying \eqref{eq:full-jet} at auxiliary precision $q$ therefore gives
\[
 |R_p|\le h\epsilon_t+(2+h)\epsilon_p+
                         \bigl(4+(d+1)h\bigr)\epsilon_q.
\]
On the unit cube $h\le1$, so the last coefficient is at most $d+5$.
Remove that auxiliary slack by \cref{lem:slack}. The resulting bound is
exactly (D), with $(2+h)\epsilon_p$. For another rational box a
primitive recursive integer upper bound on $h$ serves the same purpose.

We use the lighter notation
\begin{equation}\label{eq:jet-semantic}
 f(y)-f(x)=Df(x)(y-x)+R_f(x,y),\qquad
 |R_f(x,y)|\le h2^{-t}
\end{equation}
for the same family of rational assertions after auxiliary evaluation
errors have been eliminated by \cref{lem:slack}. This notation does not
introduce new real-number quantifiers into $\ipra$.

\subsection{Evaluation at named arguments and domain control}\label{sec:composition-eval}

A coordinatewise regular name $a$ for a point of the closed box
$D=\prod_j[\alpha_j,\beta_j]$ can be evaluated inside that box. Use
\[
 \clip_D(z)_j=\min\{\beta_j,\max\{\alpha_j,z_j\}\}.
\]
Clipping is a primitive recursive rational operation. It does not increase
the distance to any point of $D$. A point name is certified to lie in $D$
by the inequalities
$\alpha_j-\epsilon_p\le a_j(p)\le\beta_j+\epsilon_p$ for every $p$.

If $\mu_H$ is a modulus for $H$, define an extension evaluator by
\begin{equation}\label{eq:extension}
 H^\sharp(a,p)=H\bigl(\clip_D(a(q_p)),\,p+3\bigr),\qquad
 q_p=\max\{p+3,\bar\mu_H(p+3)+2\}.
\end{equation}
For $s\ge p$, the two clipped input samples differ by at most
$2^{-q_p}+2^{-q_s}\le2^{1-q_p}\le2^{-\mu_H(p+3)}$.
Applying (U) gives
\[
 |H^\sharp(a,p)-H^\sharp(a,s)|
 \le2^{-(p+3)}+2^{-(p+3)}+2^{-(s+3)}
 \le\epsilon_p+\epsilon_s.
\]
Equivalent input names produce equivalent output names by the same
argument at an increased precision. On a constant rational input this
extension agrees with the original evaluator.

For continuity assertions at named arguments one may conservatively replace
a rational modulus $\mu(t)$ by $\bar\mu(t+2)+2$. This supplies room for
the errors in both input samples. For the full-jet assertion a similar
fixed guard, for example $\bar\rho(t+2)+2$, ensures that rational
approximants lie within the required remainder radius. Approximate all
arguments, apply the rational certificate, and use \cref{lem:slack} to
remove the auxiliary errors. The sup norm, finite products, and finite sums
are continuous with the explicit error propagation given below, so this
procedure is entirely finite arithmetic at each precision.

In the remaining formulas, $\rho$ and $\mu$ refer to valid moduli for
these guarded named-argument statements. All displayed majorants remain
primitive recursive after this fixed guarding. If an implementation works
only with rational samples of an expression, the same guard can instead be
inserted during expression compilation.

This detail matters: a rational approximation to $\sigma(x)$ can fall just
outside a box that contains the exact image. Evaluating an ambient
coefficient there without a domain certificate would be unjustified.
Clipping avoids that issue. It is used only in the approximation algorithm;
we do not differentiate the clipping operation. The named composition is
the composition with $\sigma$, whose range is certified to be in $D$.

\subsection{A finite evaluator for arithmetic expressions}

Suppose represented quantities $a,b$ have known absolute bounds $A,B$.
For a sum, evaluate both at precision $p+1$. For a product, evaluate both
at
\[
 q=p+\ell(A+B+1).
\]
If $|\widetilde a-a|,|\widetilde b-b|\le2^{-q}$, then
\[
 |\widetilde a\widetilde b-ab|
 \le A2^{-q}+B2^{-q}+2^{-2q}
 \le(A+B+1)2^{-q}\le2^{-p}.
\]
Bounds for new exact values are $A+B$ and $AB$. Subtraction uses sign
reversal and addition. A finite sum uses the sum of its term bounds and
a precision shift $\ell(M)$ for $M$ terms. Each determinant is evaluated
using its finite permutation expansion.

These recursive instructions assign both a bound and an evaluation
precision to every node of a fixed arithmetic expression tree. Composition
uses \eqref{eq:extension}. Hence all expressions used in pullback and
exterior differentiation have primitive recursive evaluators with
specified error budgets. Coherence and preservation of equality follow
from the corresponding finite error inequalities and
\cref{lem:slack}.

\subsection{A local Lipschitz estimate without a mean-value theorem}

If the derivatives of a $d$-variable presentation are bounded by $D_1$,
\eqref{eq:jet-semantic} at $t=0$ gives
\begin{equation}\label{eq:local-lip}
 |f(y)-f(x)|\le(dD_1+1)h
 \qquad(h\le2^{-\rho_f(0)}).
\end{equation}
Thus
\[
 \mu_f(t)=\max\{\rho_f(0),\ t+\ell(dD_1+1)\}
\]
is an explicit value-continuity modulus. This direct estimate is enough
for every closure proof below. No existence of a maximum, mean-value
point, or integral is involved.

\subsection{Sums and products}\label{sec:product}

For $f+g$, add their jets. Use
$\rho_{f+g}(r)=\max\{\rho_f(r+1),\rho_g(r+1)\}$.
Negation preserves the modulus. The corresponding value and derivative
continuity moduli follow by the same precision shift.

For products suppose
\[
 |f|\le A_f,\quad |g|\le A_g,\quad
 |\partial_jf|\le D_f,\quad|\partial_jg|\le D_g.
\]
Put $L_f=dD_f+1$, $L_g=dD_g+1$, and take all remainder moduli to be
nondecreasing. Writing $\Delta f=f(y)-f(x)$ and similarly for $g$, direct
expansion gives
\begin{align}
 &(fg)(y)-(fg)(x)
       -\sum_j(g\,\partial_jf+f\,\partial_jg)(x)(y_j-x_j)\notag\\
 &\hspace{2em}=g(x)R_f+f(x)R_g+\Delta f\,\Delta g.
 \label{eq:product-remainder}
\end{align}
When both jets are valid at accuracy $t$ and their local Lipschitz estimates
apply, the absolute value is at most
\[
 h(A_f+A_g)2^{-t}+h^2L_fL_g.
\]
For a desired remainder accuracy $r$, set
\begin{align}
 t&=r+\ell\bigl(2(A_f+A_g+1)\bigr),\notag\\
 \rho_{fg}(r)
 &=\max\{\rho_f(t),\rho_g(t),\,r+\ell(2L_fL_g)\}.
 \label{eq:product-modulus}
\end{align}
The first term of the remainder is at most $h2^{-r}/2$, and the second
is at most $h2^{-r}/2$. Since $t\ge0$ and the moduli are monotone,
the local Lipschitz radii at accuracy zero are automatically respected.
This proves the product rule and its full uniform remainder certificate.

The product derivative is bounded by $A_gD_f+A_fD_g$. If a common
modulus for $f$ and its first derivatives is $\mu_f$, and similarly for
$g$, put
\[
 s=r+\ell(A_f+A_g+D_f+D_g+1).
\]
The modulus $\max\{\mu_f(s),\mu_g(s)\}$ controls the variation of each
product derivative: expand its two product differences, obtaining four
terms with total coefficient bound
$A_f+A_g+D_f+D_g$. This supplies the required derivative-continuity
certificate as well as the remainder certificate.

\subsection{Composition and its chain rule}\label{sec:chain-rule}

Let $f$ have $m$ inputs with $|\partial_af|\le D$, and let
$\sigma:Q^k\to D$ have componentwise $C^1$ jets with
$|\partial_i\sigma_a|\le M$. The image containment in the rational box is
part of the certificate. Let $\rho_\sigma$ majorize the component
remainder moduli and put $L=kM+1$.

For $\delta=\|y-x\|_\infty\le2^{-\rho_\sigma(0)}$,
\eqref{eq:local-lip} gives $\|\sigma(y)-\sigma(x)\|_\infty\le L\delta$.
Expand both jets:
\[
 \Delta\sigma=D\sigma(x)(y-x)+R_\sigma,\qquad
 f(\sigma(y))-f(\sigma(x))
   =Df(\sigma(x))\Delta\sigma+R_f.
\]
Subtracting the proposed chain-rule term leaves
$Df(\sigma(x))R_\sigma+R_f$. If the component remainders of $\sigma$
are at accuracy $u$, and the remainder of $f$ is at accuracy $t$, this
is bounded by
\begin{equation}\label{eq:composition-remainder}
 mD\,\delta\,2^{-u}+L\delta\,2^{-t}.
\end{equation}
For target accuracy $r$, choose
\begin{align}
 t&=r+\ell(2L),&
 u&=r+\ell\bigl(2(mD+1)\bigr),\notag\\
 \rho_{f\circ\sigma}(r)
 &=\max\{\rho_\sigma(0),\rho_\sigma(u),
                     \rho_f(t)+\ell(L)\}.
 \label{eq:composition-modulus}
\end{align}
The last condition ensures $L\delta\le2^{-\rho_f(t)}$; the other
conditions validate the two jets. Equation
\eqref{eq:composition-remainder} is at most $\delta2^{-r}$. Therefore
\begin{equation}\label{eq:chain-rule}
 \partial_i(f\circ\sigma)
 =\sum_{a=1}^m(\partial_af)\circ\sigma\,\partial_i\sigma_a
\end{equation}
with the explicit remainder modulus just given.

For completeness, derivative continuity can also be computed directly.
The derivative bound is $mDM$. In the difference of the two sums in
\eqref{eq:chain-rule}, there are $m$ product differences, with total
coefficient bound $m(M+D)$. Put
$s=r+\ell(m(M+D)+1)$. A valid modulus is the maximum of
$\rho_\sigma(0)$, the modulus for $D\sigma$ at $s$, and
the modulus for $Df$ at $s$ plus $\ell(L)$. This ensures that the
variation of every factor has the required accuracy. A value-continuity
modulus follows from \eqref{eq:local-lip} for the newly constructed jet.

The proof of \eqref{eq:chain-rule} is an algebraic expansion with finite
bounds. If named arguments occur, approximate them at an auxiliary
precision, apply these rational expansions and the expression evaluator,
and remove the final rational slack. It therefore has the same
first-order arithmetic interpretation as the cube proof.

\subsection{Second jets and symmetry}

A certified $C^2$ parameterization contains a $C^1$ jet for every
component $\sigma_a$, a $C^1$ jet for every first derivative
$\partial_i\sigma_a$, and the named equalities
\[
 H^a_{ij}=H^a_{ji},\qquad
 H^a_{ij}=\partial_i(\partial_j\sigma_a).
\]
Bounds and uniform moduli for all these finitely many fields are supplied.
Thus symmetry of the second jet is an explicit part of this presentation.
We do not silently invoke a classical theorem about interchanging
partials. Polynomial and expression-generated examples have symmetry
proofs by their finite algebraic differentiation rules. To admit a weaker
input interface containing only nonsymmetric proposed partial derivatives,
one would first need to prove the appropriate quantitative rectangular
increment theorem; that extra interface is not used here.

The following section needs exactly this much regularity: each first
derivative of $\sigma$ must have a first jet, and its antisymmetric second
part must vanish. An immersion lower bound is unnecessary for the
parameterized-chain identity.

\subsection{Polynomial certificates without integration}\label{sec:polynomial-jets}

There are many nonvacuous inputs for these interfaces. Let
$P(x)=\sum_\alpha c_\alpha x^\alpha$ be a rational polynomial on a box
$|x_j|\le R$, with integer $R\ge1$, total degree at most $D$, and integer
$A\ge\sum_\alpha|c_\alpha|$. Exact rational evaluation is primitive
recursive. Differentiation of its finite coefficient list is primitive
recursive too.

For $\|\Delta\|_\infty=h\le1$, expand each monomial by the multivariate
binomial identity. The constant and linear terms are $P(x)$ and
$DP(x)\Delta$. Every remaining term has at least two factors from
$\Delta$, hence contains $h^2$ as a bound. Summing their absolute values gives
\begin{equation}\label{eq:polynomial-jet}
 \abs{P(x+\Delta)-P(x)-DP(x)\Delta}
 \le A(R+1)^D h^2.
\end{equation}
Indeed the full binomial coefficient sum for a monomial of degree $d$
is bounded by $(R+1)^d\le(R+1)^D$.
Consequently
\[
 \rho_P(r)=r+\ell\bigl(\max\{1,A(R+1)^D\}\bigr)
\]
is a remainder modulus.

Bounds $AR^D$, $ADR^D$, and $AD^2R^D$ suffice for values, first
derivatives, and second derivatives. Continuity bounds can be derived
from the elementary product-difference identity for monomials: their
variation is at most their degree times $R$ to an appropriate power,
times $\|\Delta\|_\infty$. A common safe Lipschitz bound for the value
and its first derivatives is $A(D+1)^2R^D$. The second derivatives
are handled by differentiating the finite polynomial again. All these
claims are finite coefficient identities and inequalities; no FTC or
Stokes theorem is needed to validate them.
